\section{Literature Review}

\subsection{Transit Route Choice Models}

The theoretical foundation for route choice modeling rests on random utility
maximization \cite{benakiva1985}, under which a traveler selects the
alternative with the highest perceived utility. Prato \cite{prato2009}
identifies three persistent challenges: generating plausible alternatives,
defining attributes behaviorally, and selecting a model structure. The MNL
model \cite{mcfadden1974} remains the workhorse owing to its tractability,
though its IIA limitation has motivated richer structures. Train
\cite{train2009} demonstrated that Mixed Logit captures unobserved taste
heterogeneity through continuous coefficient distributions, while Greene and
Hensher \cite{greene2003} showed that Latent Class models segment the
population into discrete behavioral groups linked to observable
characteristics via concomitant variables. Large-scale empirical applications
remain scarce: Kim et al.\ \cite{kim2020} calibrated a transit route choice
model using Seoul smart card data but considered a single user type, Arriagada
et al.\ \cite{arriagada2025} estimated a learning-based model on
$\sim$100,000 trips in Santiago, and Arriagada et al.\ \cite{arriagada2022}
later revealed route choice strategy heterogeneity using smart card data in
the same network.

\subsection{Transfer Penalty and User Heterogeneity}

Jara-D\'iaz et al.\ \cite{jaradiaz2022}, synthesizing data from five
countries, reported pure transfer penalties of 13--18 minutes. For Seoul, Yoo
\cite{yoo2015} estimated an average penalty of 11.24 minutes but did not
disaggregate by demographics. A growing body of work documents systematic
differences across traveler groups: Chen et al.\ \cite{chen2021} found
significant preference heterogeneity among elderly riders in China, and Peng
et al.\ \cite{peng2024} reported that older rail passengers are 10.2
percentage points less likely to choose express services. These studies
examine single segments in isolation; no prior work has simultaneously
compared behavior across five user types using four model families---an
omission this paper rectifies.

\subsection{Differentiation from Prior Research}

Table~1 positions the present study relative to the most relevant prior work.
The key distinctions are the data scale (1.32 million chains), model breadth
(four families), multi-level similarity framework, and closed-loop calibration
analysis.

\begin{table}[htbp]
\centering
\tablefont
\begin{tabular}{|l|l|l|l|l|}
\hline
\headrow
\textbf{Criterion} & \textbf{\cite{kim2020}} & \textbf{\cite{yoo2015}} &
\textbf{\cite{arriagada2025}} & \textbf{This Study (2025)} \\
\hline
Data Scale    & $\sim$10K & $\sim$1K & $\sim$100K & 1,320K \\
\hline
Similarity    & Single & Single & Single & 4-Level/10 \\
\hline
Models        & MNL & MNL & MNL & 4 Families \\
\hline
User Types    & 1 & 1 & 1 & 5 + LC \\
\hline
ML Benchmark  & $\times$ & $\times$ & $\times$ & LightGBM \\
\hline
Calibration   & $\times$ & $\times$ & $\times$ & $\theta$ + Re-rank \\
\hline
\end{tabular}
\caption{Comparison with Prior Research}
\end{table}
